% ===================================================================
%  TABELUL Nr. 14 — Strada. — Negustorii.
%  Traduction 2026 d'apres texto/io, controlee sur texto/fr, texto/en
%  et texto/es. Consigne : voir texto/ro/10-tabelul-01.tex.
%
%  LE TABLEAU QUI AVAIT COUTE DEUX INVERSIONS AU SUEDOIS — le jet (49)
%  de la fontaine (48) et l'imperiale (59) de l'omnibus (58) — n'en
%  coute aucune ici : le roumain postpose son genitif comme l'ido.
% ===================================================================

%%K t14-apar-1 apar
\VUsaut{4.0mm}
\VUtitre{15.0pt}{60}{TABELUL Nr. 14}
\VUsaut{3.0mm}

%%K t14-tit-1 sub
\VUpk{11.4pt}{40}{Strada. --- Negustorii.}
\VUsaut{3.0mm}

%%K t14-01-1 p
1. --- Prietenul meu Ioan, care locuiește la țară, a venit să petreacă trei zile la
ai mei. \VUgras{Până} acum nu i s-a întâmplat niciodată să vadă un oraș mare, așa
că ne plimbăm zile întregi, iar eu îi lămuresc tot ce nu știe.

%%K t14-02-1 p
2. --- Mai întâi am mers să vedem \VUgras{observatorul} \textsuperscript{(1)}, care
l-a interesat mult. Întorcându-ne pe \VUgras{strada} \textsuperscript{(3)} unde se
află \VUgras{băile turcești} \textsuperscript{(2)}, am întâlnit o
\VUgras{înmormântare} \textsuperscript{(4)}. În fruntea \VUgras{alaiului funebru}
mergea \VUgras{clerul}, înaintea \VUgras{dricului}, pe care fusese așezat
\VUgras{sicriul}. Mulți prieteni mergeau după familia în \VUgras{doliu}.

%%K t14-02-2 p
Când alaiul a trecut, am traversat strada înaintea unei \VUgras{berării} mari
\textsuperscript{(5)}, unde mulți \VUgras{mușterii} beau \VUgras{bere} pe terasă,
sub \VUgras{șopronul} de sticlă \textsuperscript{(6)}.

%%K t14-03-1 p
3. --- Pe Ioan l-a nedumerit \VUgras{stema} pusă între prăvălia \VUgras{negustorului
de vinuri și rachiuri} \textsuperscript{(7)} și prăvălia \VUgras{vânzătorului de
note} \textsuperscript{(10)}. M-a întrebat ce înseamnă; i-am spus că este
\VUgras{consulatul} englez \textsuperscript{(8)} și i l-am arătat pe
\VUgras{consul} \textsuperscript{(9)}, care stătea pe balcon.

%%K t14-tit-2 sub
\VUpk{10.2pt}{40}{Privind vitrinele.}
\VUsaut{3.0mm}

%%K t14-04-1 p
4. --- Firește, ne-am oprit înaintea expoziției de \VUgras{instrumente muzicale},
de \VUgras{armonii}, de \VUgras{orgi} și de \VUgras{piane}; am privit copertele
ilustrate ale \VUgras{partiturilor} și ale \VUgras{bucăților} de pian și am admirat
\VUgras{lira} de lemn aurit care slujea de firmă.

%%K t14-05-1 p
5. --- Norii de praf ridicați de \VUgras{măturători} \textsuperscript{(11)} și de
\VUgras{mașina de măturat} \textsuperscript{(12)} ne-au silit să trecem mai repede
pe lângă frumoasele \VUgras{garnituri de mobilă} de la \VUgras{tâmplarul de mobilă}
\textsuperscript{(13)} și pe lângă expoziția de \VUgras{sobe} și de
\VUgras{lămpi} din \VUgras{fierărie} \textsuperscript{(14)}.

%%K t14-06-1 p
6. --- Chiar în acest loc a trebuit într-adevăr să coborâm de pe trotuar, fiindcă
era încărcat cu baloturi scoase din \VUgras{furgonul de mobilă}
\textsuperscript{(16)}, ca să fie duse în \VUgras{prăvălia} abia închiriată
\textsuperscript{(15)}.

%%K t14-06-2 p
Aceasta ne-a împiedicat să vedem trăsurile frumoase de la \VUgras{caretaș}
\textsuperscript{(17)}, dar nu ne-a împiedicat să ne oprim și să privim cum
\VUgras{ambalatorul} \textsuperscript{(19)} bate o ladă pentru vecinul său,
\VUgras{negustorul de biciclete și automobile} \textsuperscript{(18)}. Apoi, fiindcă
era deja destul de târziu, ne-am întors acasă.

%%K t14-tit-3 sub
\VUpk{10.2pt}{40}{Plimbare prin oraș.}
\VUsaut{3.0mm}

%%K t14-07-1 p
7. --- A doua zi mama a propus ca Ioan să se fotografieze și m-a rugat să îl duc la
\VUgras{fotograf} \textsuperscript{(20)}. După masă am mers în \VUgras{atelierul}
artistului, unde a trebuit să așteptăm binișor. Ca să treacă vremea, priveam
\VUgras{portretele} \textsuperscript{(22)} atârnate pe perete.

%%K t14-07-2 p
Când ne-a venit rândul, lucrul a fost scurt, și de îndată ce prietenul meu a sfârșit
de pozat înaintea \VUgras{aparatului} \textsuperscript{(21)}, am coborât.

%%K t14-08-1 p
8. --- La etajul întâi am văzut prin ușa deschisă a \VUgras{frizeriei}
\textsuperscript{(23)} cum calfa tundea un mușteriu.

%%K t14-08-2 p
Ieșind din nou în stradă, am trecut pe lângă \VUgras{tutungerie}
\textsuperscript{(24)}. Pe Ioan l-au înveselit atât de tare \VUgras{felinarul} roșu
\textsuperscript{(25)} și \VUgras{pipa mare} care slujea de \VUgras{firmă}
\textsuperscript{(26)}, încât s-a izbit de doi bătrâni care stăteau de vorbă la o
\VUgras{priză de tutun} \textsuperscript{(27)}.

%%K t14-tit-4 sub
\VUpk{10.2pt}{40}{Cofetăria.}
\VUsaut{3.0mm}

%%K t14-09-1 p
9. --- \VUgras{Bancnotele} și \VUgras{monedele} din toate țările, așezate în
\VUgras{vitrina} \VUgras{zarafului} \textsuperscript{(29)}, ne-au oprit o clipă,
dar fiindcă era ceasul gustării, am intrat fără să mai zăbovim în
\VUgras{cofetărie} \textsuperscript{(31)}.

%%K t14-10-1 p
10. --- Văzând toate \VUgras{bunătățile} care erau în prăvălie ---
\VUgras{prăjituri, turtă dulce, biscuiți} și \VUgras{bomboane} de tot felul --- cu
greu am putut alege. La urmă Ioan s-a oprit la \VUgras{prăjiturile cu cremă}, iar
eu am luat numai o \VUgras{tartă cu cireșe} mică, fiindcă de la dulciuri
\VUgras{mă dor dinții}, iar mie nu îmi place deloc să merg prea des la
\VUgras{dentist} \textsuperscript{(30)}.

%%K t14-tit-5 sub
\VUpk{10.2pt}{40}{Scrisul la mașină.}
\VUsaut{3.0mm}

%%K t14-11-1 p
11. --- Ca să îi arăt prietenului \VUgras{mașina de scris}, despre care auzise dar
pe care nu o văzuse, am intrat în \VUgras{biroul} \textsuperscript{(32)}
\VUgras{vărului} meu \textsuperscript{(34)}. L-am găsit dictând scrisori
\VUgras{secretarei} sale \textsuperscript{(35)}, care este \VUgras{stenografă}. De
îndată ce scrisoarea era gata, \VUgras{copistul} \textsuperscript{(33)} o scria la
mașina lui. Nevoind să tulburăm munca rudei, am stat la el numai câteva minute.

%%K t14-12-1 p
12. --- Sub biroul vărului locuiește un mare \VUgras{ceasornicar și bijutier}
\textsuperscript{(37)}, care pe deasupra este și argintar. În prăvălia lui minunată
sunt multe \VUgras{ceasuri de perete, ceasuri de cămin, ceasuri de buzunar},
\VUgras{lanțuri} și \VUgras{bijuterii} scumpe: \VUgras{coliere de perle},
\VUgras{brățări} de aur, \VUgras{cercei} și \VUgras{inele} cu \VUgras{pietre
scumpe} și cu \VUgras{diamante}. Una dintre vitrine este dată în întregime
\VUgras{argintului} și cuprinde o mare colecție de \VUgras{tacâmuri} de argint.

%%K t14-13-1 p
13. --- Ocolind colțul străzii, am băgat de seamă că \VUgras{tăblița}
\textsuperscript{(36)} de lângă \VUgras{numărul} casei fusese schimbată. Eram gata
să o spun cu glas tare, când s-au auzit sunetele vesele ale fanfarei militare. Am
pornit în fugă spre regiment, fără să ne gândim că va trece pe lângă prăvăliile
\VUgras{pălărierului} \textsuperscript{(39)} și \VUgras{mănușarului}
\textsuperscript{(40)} și că am fi fost tot atât de bine așezați ca să îi privim
trecerea dacă am fi rămas înaintea vitrinei \VUgras{opticianului}
\textsuperscript{(38)}, plină de \VUgras{pince-nez, ochelari} și
\VUgras{binocluri}.

%%K t14-tit-6 sub
\VUpk{10.2pt}{40}{Trecerea regimentului.}
\VUsaut{3.0mm}

%%K t14-14-1 p
14. --- În fruntea \VUgras{regimentului} \textsuperscript{(41)} mergea
\VUgras{toboșarul-șef} \textsuperscript{(42)}, după el \VUgras{toboșarii}
\textsuperscript{(43)}, \VUgras{goarnele} \textsuperscript{(44)} și toată
\VUgras{fanfara}. \VUgras{Generalul} \textsuperscript{(45)}, care era în piață cu
adjutantul său, s-a oprit lângă \VUgras{șuvoiul} \textsuperscript{(49)}
\VUgras{fântânii} \textsuperscript{(48)}, ca să privească \VUgras{parada}.

%%K t14-14-2 p
\VUgras{Ofițerii de stat-major} \textsuperscript{(46)}, \VUgras{colonelul} și
\VUgras{locotenent-colonelul} l-au salutat cu spada. \VUgras{Maiorii} mergeau în
fruntea \VUgras{batalioanelor} lor \textsuperscript{(47)}, și fiecare
\VUgras{căpitan} își comanda \VUgras{compania}, iar \VUgras{locotenenții},
\VUgras{sublocotenenții} și \VUgras{subofițerii} își țineau locurile obișnuite
lângă oamenii lor.

%%K t14-15-1 p
15. --- Mulți trecători s-au strâns la \VUgras{răspântie} \textsuperscript{(50)};
prăvăliașii --- \VUgras{umbrelarul} \textsuperscript{(51)},
\VUgras{vopsitorul} \textsuperscript{(53)}, \VUgras{armurierul}
\textsuperscript{(54)} --- au ieșit să vadă \VUgras{soldații}. Toate trăsurile ---
\VUgras{birjele} \textsuperscript{(52)}, \VUgras{trăsurile particulare}
\textsuperscript{(57)}, și de asemenea \VUgras{stropitoarea}
\textsuperscript{(56)} --- s-au lipit de trotuar.

%%K t14-tit-7 sub
\VUpk{10.2pt}{40}{Scene de pe stradă.}
\VUsaut{3.0mm}

%%K t14-15-2 p
\VUgras{Comis-voiajorul} \textsuperscript{(55)}, care își ducea în mâini
\VUgras{geamantanele cu mostre}, a trecut sprinten strada, iar cei așezați pe
\VUgras{imperiala} \textsuperscript{(59)} \VUgras{omnibuzului}
\textsuperscript{(58)} se întorceau ca să admire priveliștea. Bietul \VUgras{cântăreț
de stradă} \textsuperscript{(60)} cu \VUgras{flașneta} sa \textsuperscript{(61)} a
trebuit să își întrerupă \VUgras{cântecul}, fiindcă zgomotul tobelor îi îneca glasul.

%%K t14-16-1 p
16. --- Când regimentul a ajuns în dreptul \VUgras{prăvăliei de postavuri}
\textsuperscript{(64)}, \VUgras{croitoresele} \textsuperscript{(63)} și
\VUgras{croitorii} \textsuperscript{(62)} și-au lăsat \VUgras{mașinile de cusut} și
lucrul, ca să se uite pe fereastră. \VUgras{Prăvăliașul} \textsuperscript{(65)},
care tocmai pusese \VUgras{costume} noi \textsuperscript{(66)} în \VUgras{vitrina}
sa, a trebuit să strângă bucățile de \VUgras{postav} \textsuperscript{(67)} și de
\VUgras{pânză}, așezate pe trotuar lângă \VUgras{rigolă} \textsuperscript{(68)}, de
teamă că mulțimea le va răsturna; până și bătrânul \VUgras{negustor ambulant}
\textsuperscript{(69)}, de la care portăreasa cumpăra ciorapi, a trebuit să intre în
gang ca să sfârșească vânzarea.

%%K t14-16-2 p
Am însoțit regimentul până la cazarmă \VUgras{și}, foarte mulțumiți de ziua noastră,
ne-am întors acasă la masă.

%%K t14-tit-8 sub
\VUpk{10.2pt}{40}{Meșteșuguri și îndeletniciri felurite.}
\VUsaut{3.0mm}

%%K t14-17-1 p
17. --- A doua zi tata ne-a propus să ne ducă la tipografia ziarului local. Am
primit cu bucurie.

%%K t14-17-2 p
\VUgras{Tipograful} este totodată și \VUgras{librar} \textsuperscript{(70)},
\VUgras{editor}, \VUgras{papetar} și \VUgras{legător}. A așezat la etajul întâi
\VUgras{redacția} \textsuperscript{(71)} ziarului, și de asemenea
\VUgras{gravura}, unde lucrează \VUgras{litografii} \textsuperscript{(72)} și
\VUgras{gravorii în aramă} \textsuperscript{(73)}, și, în sfârșit, sala de cules,
unde șed \VUgras{culegătorii} \textsuperscript{(74)}. \VUgras{Sala mașinilor}
\textsuperscript{(75)} propriu-zisă, \VUgras{mașinile de tipar} și feluritele
mașini sunt la parter.

%%K t14-tit-9 sub
\VUpk{10.2pt}{40}{Ioan pune o mulțime de întrebări.}
\VUsaut{3.0mm}

%%K t14-18-1 p
18. --- Ieșind din tipografie, Ioan s-a oprit în mare nedumerire înaintea unei
cutii de sticlă de pe un stâlp din fața \VUgras{galanteriei}
\textsuperscript{(77)} și a întrebat la ce este. Tata i-a lămurit că este
\VUgras{avertizorul de incendiu} \textsuperscript{(76)}. Și într-adevăr, totul în
oraș era o minune pentru prietenul meu: și simpla \VUgras{cișmea}
\textsuperscript{(78)}, și \VUgras{tricicleta de livrare} ușoară
\textsuperscript{(80)}, pe care mergea un \VUgras{comisionar}
\textsuperscript{(79)} în livrea, și \VUgras{furgonul poștal}
\textsuperscript{(81)}.

%%K t14-19-1 p
19. --- În vreme ce tata ne-a lăsat o clipă ca să vorbească cu \VUgras{perceptorul}
\textsuperscript{(83)}, care se oprise să stea de vorbă cu \VUgras{avocatul}
\textsuperscript{(82)} și cu \VUgras{notarul} \textsuperscript{(84)}, noi ne
distram urmărind mișcarea de pe stradă. Pe lângă noi trecea cel mai felurit norod:
\VUgras{băiatul de la cofetărie} \textsuperscript{(85)}, care ducea în coș o
\VUgras{plăcintă} minunată, mergea după \VUgras{telalul} \textsuperscript{(86)};
\VUgras{lampagiul} \textsuperscript{(87)} stătea de vorbă cu \VUgras{lucrătorul de
la gaz} \textsuperscript{(88)}; o \VUgras{călugăriță} \textsuperscript{(89)}, care
ținea de mână o orfană mică, se întorcea la mănăstirea ei.

%%K t14-tit-10 sub
\VUpk{10.2pt}{40}{Pe drumul de întoarcere.}
\VUsaut{3.0mm}

%%K t14-20-1 p
20. --- Tocmai când tata s-a întors la noi, Ioan a izbucnit în râs văzând fețele
murdare și portul ciudat al doi \VUgras{coșari} \textsuperscript{(91)}, negri de
funingine.

%%K t14-21-1 p
21. --- Voiam să îi cumpăr mamei o plantă de la \VUgras{florăreasă}
\textsuperscript{(92)}, dar tata a băgat de seamă că va fi greu de dus și că este
mai bine să luăm \VUgras{portocale}. Ne-am apropiat de \VUgras{precupeața}
\textsuperscript{(94)}, iar când \VUgras{guvernanta} \textsuperscript{(93)}, care le
alegea pentru copilul ei, a sfârșit de cumpărat, am fost serviți.

%%K t14-22-1 p
22. --- Pe drumul de întoarcere am întâlnit câțiva cunoscuți: o prietenă veche a
familiei noastre, sprijinită de brațul \VUgras{doamnei ei de companie}
\textsuperscript{(95)}, un \VUgras{inginer} \textsuperscript{(96)} de poduri și
șosele și pe una dintre verișoarele mele cu \VUgras{dădaca} ei
\textsuperscript{(98)}.

%%K t14-22-2 p
Apoi am trecut pe lângă \VUgras{modista} \textsuperscript{(97)} mamei mele și pe
lângă doi \VUgras{călugări} \textsuperscript{(99)}, veniți în oraș, care s-au
apropiat să îl întrebe pe tata de drumul spre gară.
\VUsaut{6.0mm}
\VUfilet{22mm}
